% Epiplexity: definitions, estimators, follow-ups and alternatives, with consequences for small-model estimation.
% Synthesised from ../research_notes/Epiplexity literature review (notes dated 23 Sep 2026); sources in refs.bib.
% Build: pdflatex epiplexity && bibtex epiplexity && pdflatex epiplexity && pdflatex epiplexity
\documentclass[11pt]{article}
\usepackage[letterpaper,margin=1in]{geometry}
\usepackage[parfill]{parskip}
\usepackage[T1]{fontenc}
\usepackage{lmodern,microtype}
\usepackage{placeins}
\usepackage{amsmath,amssymb,amsthm,mathtools,booktabs,tabularx,array,enumitem}
\usepackage[round]{natbib}
\usepackage{xurl}
\usepackage[hidelinks]{hyperref}
\usepackage[capitalise,noabbrev]{cleveref}
\usepackage[most]{tcolorbox}

% Relaxed float fractions, as in the style model: a page may hold several tables and text together.
\renewcommand{\topfraction}{0.9}\renewcommand{\bottomfraction}{0.7}
\renewcommand{\textfraction}{0.07}\renewcommand{\floatpagefraction}{0.75}
\setcounter{topnumber}{3}\setcounter{bottomnumber}{2}\setcounter{totalnumber}{4}

% `stated': a claim the literature settles.  `unsettled': an open question, with the reason.
\newtcolorbox{stated}{enhanced jigsaw,breakable,colback=white!0,colframe=black!60,
                      left=5pt,right=5pt,top=4pt,bottom=4pt}
\newtcolorbox{unsettled}{enhanced jigsaw,breakable,colback=red!3,colframe=red!35,
                         left=5pt,right=5pt,top=4pt,bottom=4pt}

\newcommand{\E}{\mathbb{E}}\newcommand{\N}{\mathbb{N}}\newcommand{\R}{\mathbb{R}}
\newcommand{\Dt}{\mathcal{D}}\newcommand{\PT}{\mathcal{P}_T}\newcommand{\prog}{\mathrm{P}}
\newcommand{\KL}{\mathrm{KL}}\newcommand{\MDL}{\mathrm{MDL}}\newcommand{\Poly}{\mathrm{Poly}}
\newcommand{\od}{\textsuperscript{$\ast$}}
\DeclareMathOperator*{\argmin}{arg\,min}
\newtheoremstyle{tight}{4pt}{4pt}{\normalfont}{}{\bfseries}{.}{.5em}{}
\theoremstyle{tight}
\newtheorem{definition}{Definition}\newtheorem{proposition}{Proposition}
\newtheorem{remark}{Remark}
\crefname{remark}{Remark}{Remarks}\Crefname{remark}{Remark}{Remarks}
\setlist{nosep,leftmargin=*}
\setlength{\tabcolsep}{4pt}
\newcolumntype{L}{>{\raggedright\arraybackslash}X}
\newcolumntype{P}[1]{>{\raggedright\arraybackslash}p{#1}}

\title{\textbf{Epiplexity}\\[0.3em]
\large Definitions, estimators, follow-ups and alternatives,\\ with consequences for small-model estimation\\[0.3em]
\normalsize Literature to 23 September 2026}
\date{}\author{}

\begin{document}\maketitle\vspace{-3em}

\begin{abstract}\noindent
Epiplexity $S_T$ is the length of the program that minimises a two-part code for a dataset among
programs that evaluate and sample within a compute bound $T$ \citep{finzi2026epiplexity}. The
rigorous results are asymptotic and rest on cryptographic assumptions. The published values come
from two neural estimators: the area of the online loss curve above the final loss; the cumulative
teacher--student divergence of requential coding \citep{qiu2026requential}. These notes state the
definitions and theorems, compare six estimators and tabulate every trained experiment. They assess
the EpiSelect and EpiGen follow-up \citep{su2026epiplexity} and about forty papers citing the founding
one, then place older measures of structure on shared criteria. For a project with multilayer
perceptrons (MLPs) on small datasets and one Colab GPU, the main questions remain open. No published
estimate uses an MLP. No estimate reports a spread over seeds. The held-out loss floor, the units and
the per-message overhead of requential coding each shift a small-scale value by an amount comparable
to the signal.
\end{abstract}

% ============================================================================
\section{Setup and notation}\label{sec:setup}
% ============================================================================
Logarithms are base 2 and code lengths are in bits. The released code of every paper below logs
natural-log losses; a value in nats becomes bits after division by $\ln 2$
\citep{epiplexitycode,requentialcode}. $\KL(p\|q)$ is the Kullback--Leibler divergence in bits.

$X$ is the dataset whose information is measured, of $\Dt$ tokens or examples;
\citet{finzi2026epiplexity} call it the test dataset. A training run sees units $Z_0,\dots,Z_{M-1}$
in order; $P_i$ is the model after $i$ updates. When every unit is a token, $M=D$, the number of
training tokens. $N$ is the number of parameters, embeddings included in the reference code
\citep{epiplexitycode}. Compute is measured in floating-point operations (FLOPs). Training for $D$
tokens takes $6ND$ FLOPs and evaluating on $X$ takes $2N\Dt$ \citep{finzi2026epiplexity}. A run is
admissible under the compute bound $T$ when
\begin{equation}\label{eq:T}
6ND+2N\Dt\le T .
\end{equation}
One Colab GPU means one T4, L4 or A100 in a single session.

``(Our derivation)'' marks an inference absent from the cited source; ``(our measurement)'' marks a
value obtained by running released code on a CPU; in tables \od\ marks either. The sign $\approx$
marks a value estimated from a published figure, with an error of roughly ten per cent.

% ============================================================================
\section{Definitions and theorems}\label{sec:defs}
% ============================================================================
\begin{definition}[$T$-time probabilistic model; Def.~7 of \citealp{finzi2026epiplexity}]\label{def:model}
Let $T:\N\to\N$ be non-decreasing and time-constructible. Let $\mathcal U$ be a fixed prefix-free
universal Turing machine. A prefix-free program $\prog$ is a $T$-time probabilistic model over
$\{0,1\}^n$ if both modes of $\prog$ halt within $T(n)$ steps: $\mathcal U(\prog,(0,x))$ outputs a
probability $P(x)$ with a finite binary expansion; $\mathcal U(\prog,(1,u))$ outputs a sample in
$\{0,1\}^n$ from an infinite random tape $u$. The probabilities sum to one and the samples follow
$P$. $\PT$ is the set of such programs; $|\prog|$ is the length of $\prog$ in bits.
\end{definition}

$K(x)$ denotes the prefix Kolmogorov complexity, the length of the shortest program for which
$\mathcal U$ outputs $x$ \citep{livitanyi2019}.

\begin{definition}[epiplexity; Def.~8 of \citealp{finzi2026epiplexity}]\label{def:epi}
For a random variable $X$ on $\{0,1\}^n$, let
\[
\prog^\star=\argmin_{\prog\in\PT}\bigl\{|\prog|+\E[\log 1/P(X)]\bigr\},
\]
with ties broken by the shortest program. The epiplexity is $S_T(X):=|\prog^\star|$. The
time-bounded entropy is $H_T(X):=\E[\log 1/P^\star(X)]$. The sum $\MDL_T(X):=S_T(X)+H_T(X)$ is the
total time-bounded information, a minimum description length (MDL).
\end{definition}

\begin{definition}[conditional; Def.~11 of \citealp{finzi2026epiplexity}]\label{def:cond}
$S_T(Y\mid X)$ and $H_T(Y\mid X)$ follow from \cref{def:epi} with $\E[\log1/P(X)]$ replaced by
$\E_{(X,Y)}[-\log P(Y\mid X)]$ and the minimum taken over programs whose conditional $P_{Y\mid x}$
belongs to $\PT$ for every fixed $x$.
\end{definition}

In general $H_T(Y,X)-H_T(X)\neq H_T(Y\mid X)$. A conditional estimate sums code length over label
tokens only, while the time bound includes input tokens \citep[App.~B.1]{finzi2026epiplexity}. A
classifier therefore estimates $S_T(Y\mid X)$ (our derivation). Every estimator below replaces $\PT$
by the models one optimiser reaches from one architecture, a restriction the authors allow in a
footnote to Def.~7.

\begin{proposition}[basic properties; Sec.~3 of \citealp{finzi2026epiplexity}]\label{prop:basic}
Let $c_1,c_2$ be constants independent of $X$ and $H(X)$ the Shannon entropy.
(i)~$S_T,H_T\ge0$.
(ii)~$H(X)\le\MDL_T(X)\le n+c_1$.
(iii)~$\MDL_{T'}\le\MDL_T$ whenever $T'\ge T$.
(iv)~$\MDL_{T'}(f^{-1}(X))\le\MDL_T(X)+|\mathrm f|+c_2$ with $T'=T+\mathrm{Time}(\mathrm f)$, for a
program $\mathrm f$ computing a bijection in fixed time.
(v)~For uniform $U_n$ and any $T(n)\ge c_1$, $S_T(U_n)\le c_2$. The equal mixture of $0101\ldots$
and $1010\ldots$ has $S_T=O(1)$ and $H_T=O(1)$ under a linear time bound.
\end{proposition}

Only $\MDL_T$ is monotone in $T$. $S_T$ can fall as $T$ grows, as in the emergence experiment of
\cref{tab:exp}; a study should report the curve $T\mapsto S_T$ (our derivation).

\Cref{tab:thm} collects the theorems. A pseudorandom generator (PRG) $G$ stretches $k$ uniform bits
$U_k$ to $n=\mathrm{poly}(k)$ bits. No non-uniform probabilistic polynomial-time (PPT) test separates
$G(U_k)$ from $U_n$ with advantage above $\varepsilon(k)$. A one-way function (OWF) is computable in
polynomial time and hard to invert for such tests. The subscript $\Poly$ means that the bound holds
for every polynomial time bound $T(n)$.

\begin{table}[!htb]
\centering\small
\caption{Theorems of \citet{finzi2026epiplexity}, v2 numbering; $c$ is a constant. Paradox~3 has a
definition and a conjecture but no theorem.}\label{tab:thm}
\begin{tabularx}{\textwidth}{@{}P{2.3cm}P{3.6cm}L@{}}
\toprule
Result & Assumption & Statement \\
\midrule
Thm~9 (PRG) & PRG $G$ with advantage $\varepsilon(k)$ &
$n-2-n\varepsilon(k)<H_\Poly(G(U_k))\le n+c$; $S_\Poly(G(U_k))\le c+n\varepsilon(k)$; Shannon entropy $k$ \\
Thm~10 & OWFs secure against non-uniform PPT & some $X_n$ has $S_\Poly(X_n)=\Omega(\log n)$; non-constructive \\
Thm~12 (Paradox~1) & as Thm~9 & $H_\Poly(G(U_k))-H_\Poly(U_k)>n-k-n\varepsilon(k)-c$ \\
Thm~13 (Paradox~2) & one-way permutation $f$; $X=U_n$; $Y=f(X)$ & $H_\Poly(X\mid Y)+H_\Poly(Y)>H_\Poly(Y\mid X)+H_\Poly(X)+\omega(\log n)$; by Cor.~26 no polynomial-time model fitting the forward direction obeys Bayes' rule \\
Def.~14 (Paradox~3) & none & emergence: $S_{T_1}-S_{T_2}=\Theta(1)$ for one step, $\omega(1)$ for $k(n)$ steps, $T_1=o(T_2)$; no instance proven \\
Thm~30; App.~B.4 & loss $C^2$ with a unique optimum, larger models more sample-efficient & optimal $N^\star(T)$ and $D^\star(T)$ increase with $T$; $D^\star\to\Dt$ from below and $N^\star\sim T/(8\Dt)$; with infinite compute $S_\infty$ is non-decreasing in $\Dt$ \\
\bottomrule
\end{tabularx}
\end{table}

\begin{stated}
Under cryptographic assumptions and polynomial time bounds the theory settles three claims.
Pseudorandom data has near-maximal time-bounded entropy and near-constant epiplexity. A deterministic
map can create time-bounded entropy. Time-bounded entropy depends on the order of factorisation. No
theorem covers the finite, FLOP-bounded estimators of \cref{sec:est}; those rest on the scaling
assumptions of App.~B.4 of \citet{finzi2026epiplexity}.
\end{stated}

\begin{unsettled}
Does any explicit distribution have epiplexity above $\Omega(\log n)$? Does computation create
epiplexity, as opposed to entropy? Reason: the existence proof is non-constructive; Theorem~12
concerns $H_T$; the $S_T$ claims of Paradoxes~1 and~3 rest on the cellular-automaton and looped-model
experiments; no instance of Definition~14 is proven \citep{finzi2026epiplexity}.
\end{unsettled}

% ============================================================================
\section{Estimators}\label{sec:est}
% ============================================================================
\subsection{Prequential}\label{sec:preq}
Prequential coding \citep{dawid1984} sends $Z_i$ with $\log1/P_i(Z_i)$ bits and then trains on it.
The total encodes the data together with the final model and decodes in $6ND$ FLOPs. Following
\citet{zhang2020transfer} and \citet{finzi2025computeoptimal}, Eq.~8 of \citet{finzi2026epiplexity}
estimates the model part as
\begin{equation}\label{eq:preq}
|\prog_{\rm preq}|\approx\sum_{i=0}^{M-1}\bigl[\log1/P_i(Z_i)-\log1/P_M(Z_i)\bigr].
\end{equation}
The right side is the area of the online loss curve above the final loss. The two-part code adds
$\E[\log1/P_M(X)]$, estimated as the held-out loss per token times $\Dt$. The authors give two reasons
why \eqref{eq:preq} is not rigorous. The difference of two upper bounds on Kolmogorov complexities does
not bound $K(P_M)$. Symmetry of information fails for time-bounded complexity, which leaves the
decoding time unproven.

Every experiment except the comparison with adaptive data optimisation (ADO; \citealp{jiang2025ado})
replaces $\sum_i\log1/P_M(Z_i)$ by $M$ times the loss of
$P_M$ on unseen data \citep[App.~B.1]{finzi2026epiplexity}. The replacement assumes one pass over
i.i.d.\ data with a small generalisation gap. \citet{zhang2020transfer} compute the same difference on
examples outside the training set (information transfer); on CIFAR-10 with random labels it remains
near zero although training accuracy is high. Let $p_{\rm data}$ be the data distribution. By Eq.~1
of \citet{qiu2026requential}, the expectation of \eqref{eq:preq} with a held-out floor is (our
derivation)
\begin{equation}\label{eq:preqkl}
\E|\prog_{\rm preq}|=\sum_{i=0}^{M-1}\bigl[\KL(p_{\rm data}\|P_i)-\KL(p_{\rm data}\|P_M)\bigr].
\end{equation}
The data entropy cancels. The raw prequential code $\sum_i\log1/P_i(Z_i)$ keeps it and grows linearly
after learning stops \citep{qiu2026requential}.

\subsection{Requential}\label{sec:req}
Requential coding transmits a student model \citep{qiu2026requential}. Encoder and decoder share the
student initialisation, update rule, pseudorandom seed, number of steps and batch size. At step $i$
both sides generate proposals from the student $P^{\rm s}_i$. The encoder holds a teacher $P^{\rm t}_i$
trained on real data. By relative entropy coding (REC) it selects a proposal distributed as $P^{\rm
t}_i$ and sends the index with an Elias-delta code. Both sides then train the student on that batch.
Only hard samples are coded. With $\KL_i=\KL(P^{\rm t}_i\|P^{\rm s}_i)$ and the sum over student
steps, Eq.~2 of \citet{qiu2026requential} bounds the expected length:
\begin{equation}\label{eq:req}
|\prog_{\rm req}|\le\sum_i\bigl[\KL_i+2\log(1+\KL_i)+\kappa\bigr],\qquad\kappa<5.21 .
\end{equation}
Eq.~9 of \citet{finzi2026epiplexity} has $\log(1+\KL_i)+4$ in place of the last two terms; the factor 2
comes from coding the index with a universal integer code. Both papers report $\sum_i\KL_i$. The
remaining terms add 5 to 15 bits per step and are negligible only when every $\KL_i$ is far above that
(our derivation).

Measuring the length needs no REC search: sampling from the teacher gives the divergences at about
$7/3$ of the FLOPs of training the teacher alone. Transmission would need about $2^{\KL_i}$ proposals
per message; no encoder is released \citep{qiu2026requential}.

The two papers keep $\KL_i$ small in different ways. The founding paper distils from an exponential
moving average (EMA) of the teacher and freezes the teacher while $\KL_i$ exceeds a threshold. The
later paper keeps the EMA and adds iso-loss projection: at steps $100\cdot1.5^j$ the teacher is reset
to the student and retrained on real data, with the student paused, until it recovers the earlier
validation loss. With a static teacher close to $p_{\rm data}$, requential coding reduces to
\eqref{eq:preq}, which the founding authors therefore expect to exceed the requential code
\citep[App.~B.2]{finzi2026epiplexity}.

For a classifier over $C$ classes with batch inputs $x_b$ known to the decoder,
$\KL_i=\sum_b\sum_{c=1}^{C}P^{\rm t}_i(c\mid x_b)\log[P^{\rm t}_i(c\mid x_b)/P^{\rm s}_i(c\mid x_b)]$
is exact and needs one forward pass of each network. Sampling a label replaces autoregressive
generation (our derivation).

\subsection{From a code length to \texorpdfstring{$S_T$}{S\_T}}\label{sec:sweep}
\citet[App.~B.1]{finzi2026epiplexity} obtain $S_T$ from a sweep. They tune the learning rate on a
small model and transfer it by $\mu$P and CompleteP. They train a grid of widths and depths, each for
more tokens than $\Dt$, with a constant learning rate and an EMA of the iterates; every logged
checkpoint becomes a point (compute of \eqref{eq:T}, $|\prog|+\E\log1/P(X)$). $S_T$ is the model part
at compute $T$ on the lower convex hull of all points, of which only the median point per run is kept. The
grid has 30 configurations per cellular-automaton rule and 48 per natural dataset, each with one seed
\citep{epiplexitycode}. When a model reaches a known loss floor, as in the induction tasks, one model
per condition suffices. For ranking datasets a single run of \eqref{eq:preq} is the crude estimate the
authors recommend. They expect the hull heuristic, the fixed architecture and imperfect
hyperparameters to cause only sub-leading errors.

\begin{proposition}[power-law loss; App.~B.3 of \citealp{finzi2026epiplexity}]\label{prop:pl}
Let the loss after $D$ tokens with $N$ parameters be $\mathcal L(N,D)=E+(N_0/N)^\alpha+(D_0/D)^\beta$
with $0<\alpha,\beta<1$. Write $\hat N=N/N_0$, $\hat D=D/D_0$, $\hat\Dt=\Dt/D_0$ and
$\hat T=T/(2N_0D_0)$; equality in \eqref{eq:T} becomes $\hat T=\hat N(3\hat D+\hat\Dt)$. Take
\eqref{eq:preq} as model part and $\Dt\,\mathcal L$ as data part. Then:
\begin{enumerate}
\item[(a)] $|\prog_{\rm preq}|=\frac{\beta}{1-\beta}D_0\hat D^{1-\beta}$ for large $D$;
\item[(b)] the optimum solves $\beta\hat D^{-\beta-1}(\hat\Dt-\hat D)=3\alpha\hat\Dt\,\hat
T^{-\alpha}(3\hat D+\hat\Dt)^{\alpha-1}$ with $\hat N^\star=\hat T/(3\hat D^\star+\hat\Dt)$;
\item[(c)] as $T\to\infty$, $\hat D^\star\to\hat\Dt$ and $\hat N^\star\sim\hat T/(4\hat\Dt)$, with limits
$S_\infty(X)=\frac{\beta}{1-\beta}D_0^{\beta}\Dt^{1-\beta}$ and $H_\infty(X)=\Dt E+D_0^\beta\Dt^{1-\beta}$;
\item[(d)] for $\hat D^\star\ll\hat\Dt$, $\hat N^\star\approx\hat T/\hat\Dt$, $S_T\propto
T^{\alpha(1-\beta)/(1+\beta)}$ and $H_T-\Dt E\propto T^{-\alpha\beta/(1+\beta)}$.
\end{enumerate}
\end{proposition}
Code lengths are in the units of $E$, nats in the fitted laws. With the Chinchilla exponents
$\alpha\approx0.34$, $\beta\approx0.28$ the exponents in (d) are $0.19$ and $-0.07$. For language
($\beta=0.285$, $D_0=1.49\times10^9$, $\Dt=10^{12}$), part (c) gives $9.0\times10^{10}$ bits, the value
of Fig.~8b there (our derivation).

\subsection{Alternative estimators from 2026}\label{sec:alt}
\paragraph{Excess description length (EDL).} \citet{donoway2026edl} posted EDL two days after
epiplexity; a single-author version appeared at ISIT 2026 \citep{donoway2026edlisit}. Let
$\ell(\theta;x,y)=-\log p_\theta(y\mid x)$, $L(\theta)=\E_{p_{\rm data}}\ell$ the population loss,
$\mathcal A$ the training algorithm, $\theta_0$ the initial parameters and $\theta_T$ the parameters
after compute $T$ (steps, epochs or FLOPs). For a sample $S=(z_1,\dots,z_M)$ the ISIT definition is
\begin{equation}\label{eq:edl}
\mathrm{EDL}_T(\mathcal A,M)=\E_{S\sim p_{\rm data}^M}\Bigl[\E_\pi\sum_{i=1}^{M}
\ell\bigl(\theta^\pi_{i-1};z_{\pi(i)}\bigr)-M\,L(\theta_T)\Bigr].
\end{equation}
Here $\pi$ is a uniform permutation and $\theta^\pi_{i-1}$ are the parameters after $i-1$ examples of
the first pass; later epochs enter only through $L(\theta_T)$. The ISIT version proves non-negativity,
monotonicity in compute, additivity, saturation at $M$ times the input--label mutual information
$I(X;Y)$ and $\mathrm{EDL}_T\lesssim0$ for random labels. Neither version reports a neural-network experiment or code. EDL evaluates a given
algorithm, subtracts a population loss and allows compute beyond one pass
\citep{donoway2026edlisit}. For one pass over i.i.d.\ data it equals the expectation of
\eqref{eq:preq} with the population loss as floor (our derivation). The supremum over algorithms is
not a structure measure: an algorithm idle during the first pass reaches
$M\,[L(\theta_0)-L(\theta_T)]$, linear in $M$ even for a one-line rule (our derivation).

\paragraph{Reservoir score.} \citet{zhang2026learnablenovelty} replace training by a closed form. A
frozen random network $\phi$ maps $N_{\rm s}$ inputs to $m$ standardised features $\tilde
H\in\R^{N_{\rm s}\times m}$; $\tilde Y$ are the targets, centred and divided by a fixed scale. A
ridge-regression head with penalty $\lambda$, $W_\lambda=(\tilde H^\top\tilde H+\lambda I_m)^{-1}\tilde
H^\top\tilde Y$, with singular values $s_i$ and a resolution $\eta$ (1, or 30 for MNIST) gives
\begin{equation}\label{eq:zl}
S^\phi(Y\mid X)=\tfrac12\sum_i\log_2\bigl(1+\eta\,s_i^2\bigr).
\end{equation}
Ridge approximates the two-part minimiser; the exact minimiser ranks automaton rules with Spearman
correlation 0.997. The score is capacity-bounded, not time-bounded, with the prefactor fixed at
$\frac12$ by convention; the bits of $S^\phi$ are not commensurable with $S_T$ (our derivation). On elementary
cellular automata (ECA) it ranks rule 110 first and rule 54 above rule 30, an order that reverses once
the receptive field exceeds radius two. The authors claim that noise contributes nothing to the model
part. \Cref{tab:zl} shows a positive floor at finite $N_{\rm s}$ instead: independent targets score
close to rule 15 and decay slowly with $N_{\rm s}$ (our measurement). Rule 30 scores above rule 15,
opposite to \citet{finzi2026epiplexity}; both exceed the noise floor by a few bits only. Comparisons
need a shuffled-target baseline at matched $N_{\rm s}$. Three independent repositories reused the score within two months
\citep{lcrhepiplexity,epijepa,eberlful}.

\paragraph{Small-data prequential estimate.} \citet{liu2026selfplay} adapt \eqref{eq:preq} to
multi-epoch fine-tuning on small sets \citep{selfplaycode}. The online loss $L_{\rm online}$ over the
$N_{\rm train}$ tokens of the first epoch is frozen. After each epoch the model is evaluated on all
training examples ($L_{\rm train}$) and on a validation split ($L_{\rm val}$ over $N_{\rm val}$ tokens).
The estimate is $S=(L_{\rm online}-L_{\rm train})/\ln2$ at the epoch minimising $S/N_{\rm train}+L_{\rm
val}/(N_{\rm val}\ln2)$. With $M=N_{\rm train}$ units, $L_{\rm train}$ is in-sample, which gives
$S_{\rm Liu}=\mathrm{EDL}_T+M\,[L(\theta_T)-L_{\rm train}/M]$ in expectation (our derivation).
Memorisation raises it; only the epoch choice limits it. The code
evaluates online losses with dropout active and final losses without it, which biases $S$ upward (our
derivation from the code).

\paragraph{Finite-time thermodynamics.} \citet{ohzeki2026nonequilibrium} extends the widely
applicable Bayesian information criterion to finite-time stochastic-gradient Langevin ensembles. The
result is a computable thermodynamic term of time-bounded MDL, which the author states is not
epiplexity itself. The paper is theoretical.

\Cref{tab:est} compares the six estimators on the properties that define epiplexity.

\begin{table}[!htb]
\centering\footnotesize
\caption{Estimators of learnable structure. Entries without \od\ are claims of the cited papers.}\label{tab:est}
\begin{tabularx}{\textwidth}{@{}P{2.9cm}P{2.1cm}P{1.6cm}P{1.4cm}P{1.6cm}LP{1.7cm}@{}}
\toprule
Estimator & Zero on noise & Zero on trivial data & Compute-bounded & Minimises over models & Training runs & Code \\
\midrule
Prequential \eqref{eq:preq}; held-out variant \citep{zhang2020transfer} & with a held-out or one-pass floor\od & yes & yes & sweep, hull & 5--30 per grid of $T$; 1 to rank & PyTorch, JAX \citep{epiplexitycode} \\
Requential \eqref{eq:req} & yes, \cref{tab:req} & yes, \cref{tab:req} & yes & founding paper only & $\approx7/3$ of a run per configuration; 2--10 times prequential & JAX \citep{requentialcode} \\
EDL \eqref{eq:edl} & $\lesssim0$, proved & zero with calibrated start & indexed by $T$ & no & 1 plus held-out loss & none \\
Reservoir $S^\phi$ \eqref{eq:zl} & no at finite $N_{\rm s}$\od & zero on constants\od & no & no & none; $<1$\,s on a CPU\od & PyTorch \citep{learnablenoveltycode} \\
$S_{\rm Liu}$ \citep{liu2026selfplay} & via epoch choice\od & yes\od & no & epochs only & 1 multi-epoch run & PyTorch \citep{selfplaycode} \\
Free energy \citep{ohzeki2026nonequilibrium} & not stated & not stated & finite time & not stated & Langevin ensemble & none \\
\bottomrule
\end{tabularx}
\end{table}

\begin{table}[!htb]
\centering\small
\caption{Reservoir score $S^\phi$ in bits (our measurement with the released code of
\citealp{learnablenoveltycode}, CPU). Configuration: depth 3, 256 channels, kernel 3,
$\lambda=0.03$, $\eta=1$, target window 32. Noise: targets independent of inputs. One reservoir seed
per row except the last, which gives mean $\pm$ s.d.\ over three seeds.}\label{tab:zl}
\begin{tabular}{@{}lrrrrr@{}}
\toprule
$N_{\rm s}$ & rule 110 & rule 54 & rule 30 & rule 15 & noise \\
\midrule
128  & 62.4 & 38.7 & 34.4 & 29.9 & 28.4 \\
512  & 63.2 & 36.6 & 24.9 & 20.3 & 17.6 \\
2048 & 63.7 & 35.5 & 16.8 & 11.0 & 9.5 \\
512, three seeds & $63.9\pm0.6$ & $37.1\pm0.4$ & $25.5\pm1.0$ & $20.9\pm0.6$ & $18.0\pm1.3$ \\
\bottomrule
\end{tabular}
\end{table}

% ============================================================================
\section{Experiments of the founding and requential papers}\label{sec:exp}
% ============================================================================
\Cref{tab:exp} lists every trained experiment of \citet{finzi2026epiplexity} and
\citet{qiu2026requential}. All use GPT-2-style transformers trained with Adam and $\mu$P
learning-rate transfer; neither paper trains an MLP or a convolutional network. The founding paper
used six 2080 Ti GPUs for small experiments, six Titan RTX GPUs and 32 TPU v4 chips for natural data.
The requential paper used TPU v6e-8 hosts. Neither reports GPU-hours. Every sweep except the
one-way-function experiment (two seeds) uses one seed; no figure has error bars
\citep{epiplexitycode,requentialcode}.

\begin{table}[!htb]
\centering\footnotesize
\caption{Trained experiments. $\approx$: estimated from figures. FLOPs with \od: $6ND$ computed by us.
Last column: our estimate at 5--20 TFLOP/s sustained, not a measurement.}\label{tab:exp}
\begin{tabularx}{\textwidth}{@{}P{3.1cm}P{2.4cm}P{1.9cm}P{1.7cm}LP{1.2cm}@{}}
\toprule
Experiment & Models & Data & FLOPs & Main result & Colab \\
\midrule
\multicolumn{6}{@{}l}{\emph{\citet{finzi2026epiplexity}}}\\
ECA rules 15, 30, 54; $Y=F^{48}(X)$, 64 cells (Fig.~3) & width 16--512, depth 1--9 & $\Dt=10^8$ label tokens & $\approx10^{11}$--$2\times10^{17}$ & $S_T$: rule 54 $\approx5.4\times10^6$ bits, rule 15 $\approx2\times10^5$, rule 30 $\approx0$; order settles from $\approx10^{13}$ FLOPs & width $\le128$ \\
Ten ECA rules (Fig.~2c) & width 16--128, depth 1--3 & $\Dt=2.5\times10^8$ & $\le10^4$ steps & prequential several times requential; correlated within groups & yes \\
Rule 30 as one-way function (Fig.~4a) & 8 layers, width 128--192 & $n=16$--64 cells & $\approx10^{16}$--$10^{17}$ & $H_T$ only: reverse gap grows with $n$ & $n\le32$ \\
Chess move order; transfer (Figs.~4c, 7) & $10^6$--$1.6\times10^8$ & $5\times10^9$ chars & $\le5\times10^{18}$ & reverse order: higher $H_T$ and $S_T$ ($\approx2.8$ against $2.4\times10^8$ bits); centipawn accuracy $\approx0.28$ against $0.21$ & no \\
Hard induction, $h$ hidden bits (Fig.~5b) & 3 layers, width 256 & $\approx2\times10^9$ tokens & $\le2\times10^{16}$ & $S$ rises with $h$, $\approx2.5$ to $5.5\times10^6$ (nats) & 1--3 hours each \\
Easy induction, Markov chains (Fig.~5c) & 3 layers, width 128 & $\approx6\times10^8$ tokens & $\le2.5\times10^{15}$ & $S$ peaks at intermediate $h$: $\approx10^6$ at $h=0$, $4.7\times10^6$ at $h=6$ & minutes \\
Emergence, rule 54, looped (Fig.~6) & width 16--128, depth 1--32, 1--16 loops & $10^9$ tokens per run & $\le4\times10^{16}$ per run\od & $S_T$ and $\MDL_T$ drop abruptly when the looped model wins; peak $\approx1.45\times10^7$ bits & subset \\
Natural data (Fig.~8a) & $\le1.6\times10^8$ & $5\times10^9$ tokens & $\le6\times10^{18}$ & $S_T$: text $\approx3\times10^8$, chess $\approx2$--$2.8\times10^8$, CIFAR-5M $\approx9\times10^7$ bits & no \\
Scaling laws (Fig.~8b) & published fits & $\Dt=10^{12}$ & $T=10^{25}$ & language above images above video & CPU \\
ADO sampling (Fig.~8c) & $1.3\times10^9$ & $1.25\times10^{11}$ & $\approx10^{21}$\od & ADO $\approx2.45$ against $1.95\times10^{10}$ bits & no \\
\midrule
\multicolumn{6}{@{}l}{\emph{\citet{qiu2026requential}}}\\
Codes compared; learnable information (Figs.~4, 9) & $\approx10^8$ & $2$--$20\times10^9$ & $\approx10^{18}$--$10^{19}$\od & \cref{tab:req} & toy runs \\
Scaling, bounds, many epochs (Figs.~5, 7, 8) & $1.7\times10^6$--$10^9$ & $10^8$--$2\times10^{10}$ & $\le10^{20}$\od & code for a fixed loss falls with $N$; bound below 4-bit quantisation; best bound near one epoch & no \\
\bottomrule
\end{tabularx}
\end{table}

\begin{table}[!htb]
\centering\small
\caption{Code lengths in bits digitised from Figs.~4 and~9 of \citet{qiu2026requential}; about
$10^8$ parameters. Heuristic: \eqref{eq:preq}. Best: EMA smoothing and iso-loss projection.}\label{tab:req}
\begin{tabular}{@{}llrrrr@{}}
\toprule
Fig. & Data & Prequential, raw & Heuristic & Requential, plain & Requential, best \\
\midrule
4 & OpenWebText (characters) & $3.8\times10^9$ & $9.0\times10^8$ & $1.0\times10^9$ & $7.0\times10^7$ \\
4 & CIFAR-5M (pixels) & $9.8\times10^9$ & $2.3\times10^8$ & $4.9\times10^8$ & $3.3\times10^7$ \\
4 & FineWeb (BPE) & $1.3\times10^{10}$ & $1.75\times10^9$ & $1.7\times10^9$ & $3.2\times10^8$ \\
9 & uniform random tokens & $3.27\times10^{10}$ & -- & -- & $1.4\times10^5$ \\
9 & one repeated token & $1.7\times10^8$ & -- & -- & $1.0\times10^6$ \\
9 & CIFAR-5M & $2.3\times10^{10}$ & -- & -- & $3.6\times10^7$ \\
9 & FineWeb & $2.8\times10^{10}$ & -- & -- & $3.6\times10^8$ \\
\bottomrule
\end{tabular}
\end{table}

The ECA frontier is the smallest demonstration of the ordering structured above trivial above
pseudorandom. Order dependence of $S_T$ is measured only on chess; the small rule-30 experiment reports
$H_T$ only. In \cref{tab:req} the raw prequential code ranks uniform noise highest and the requential
code ranks it lowest. The ratio of heuristic to requential code changes several-fold with the teacher
recipe, which makes it a property of the recipe as much as of the data (our derivation).

% ============================================================================
\section{EpiSelect and EpiGen}\label{sec:follow}
% ============================================================================
\citet{su2026epiplexity} present both methods in one paper by the same group (arXiv v1, 12 August
2026). Both replace the final-loss floor of \eqref{eq:preq} by the current loss, which makes the
quantity optimisable online. Neither searches over $(N,D)$ at fixed $T$; the authors call the result
``a tractable proxy'' and leave the link to $S_T$ open.

EpiSelect reweights data domains during pretraining. Let $L_j(s)$ be the training loss of domain
$j$ at step $s$ and $D_k$ the number of tokens seen from domain $k$. The online proxy, a cross-domain
scaling law, the gradient and the sampling rule are Eqs.~2--5 there:
\begin{gather}
\hat S(t)=\sum_{j}\sum_{s=1}^{t}\bigl[L_j(s)-L_j(t)\bigr],\qquad
\hat L_j=\epsilon_j+b_j\Bigl(\sum_k\gamma_{jk}D_k\Bigr)^{-a_j},\label{eq:episelect}\\
\frac{\partial\hat S}{\partial D_k}=\sum_j\frac{a_j\gamma_{jk}D_j}{\sum_i\gamma_{ji}D_i}\bigl(\hat
L_j-\epsilon_j\bigr),\qquad \pi_k\propto\exp\Bigl(\frac1\tau\,\frac{\partial\hat S}{\partial
D_k}\Bigr).\notag
\end{gather}
The weights $\gamma_{jk}>0$ with $\sum_k\gamma_{jk}=1$ couple domains, $\epsilon_j$ is the irreducible
loss, $\pi_k$ is the sampling probability of domain $k$ and $\tau=1$. EpiSelect refits the law every
1000 steps from 729 L-BFGS initialisations and averages weights with momentum 0.9
\citep{episelectcode}. With $\gamma$ the identity the gradient reduces to $a_k(\hat L_k-\epsilon_k)$:
the ADO slope \citep{jiang2025ado} multiplied by $D_k$, without prior and credit terms (our
derivation).

EpiGen fine-tunes a GPT-2 generator by REINFORCE. The reward for a generated batch is the drop in
learner loss on a buffer $\mathcal B$ of past batches after 10 learner steps on the new batch. With
$\theta_t$ the learner parameters after update $t$, Eq.~6 there is
\begin{equation}\label{eq:epigen}
r_t=\sum_{x\in\mathcal B}\bigl[\log1/P_{\theta_{t-1}}(x)-\log1/P_{\theta_t}(x)\bigr].
\end{equation}
By design, memorising one noise batch does not lower loss on the rest of the buffer. \Cref{tab:follow}
gives the results.

\begin{table}[!htb]
\centering\small
\caption{EpiSelect: mean zero-shot accuracy over 10 LM-Eval tasks after 15.7B Common Pile tokens.
EpiGen: mean over 7 GLUE tasks and OpenWebText perplexity. From Tables~1--3 of
\citet{su2026epiplexity}; one seed per entry. Columns with \od\ drop the two tasks with the largest
gains (our computation).}\label{tab:follow}
\begin{tabular}{@{}lrrrr@{}}
\toprule
EpiSelect & 124M & 1.3B & 124M without BBQ, BoolQ\od & 1.3B without BBQ, BoolQ\od \\
\midrule
Natural sampling & .377 & .422 & .362 & .413 \\
ADO & .379 & .425 & .367 & .417 \\
EpiSelect & .394 & .431 & .368 & .420 \\
\bottomrule
\end{tabular}\\[6pt]
\begin{tabular}{@{}lrrr@{}}
\toprule
EpiGen (GPT-2 small) & GLUE & GLUE without CoLA, RTE\od & perplexity \\
\midrule
Pretrained & .743 & .860 & 24.64 \\
Frozen generator & .759 & -- & 26.49 \\
Perplexity reward & .756 & -- & 104.20 \\
No buffer & .757 & -- & 33.10 \\
EpiGen & .770 & .858 & 27.58 \\
\bottomrule
\end{tabular}
\end{table}

The evidence has four limits.
\begin{enumerate}
\item Every entry of \cref{tab:follow} comes from one seed \citep{episelectcode,epigencode}. The gains
concentrate in two tasks per method; without them EpiSelect and ADO differ by about a tenth of a point
at 124M, while EpiGen no longer exceeds the pretrained model (our computation).
\item The motivating correlation between prequential epiplexity and mean out-of-distribution (OOD)
accuracy (Pearson 0.88,
$p\approx0.05$) rests on five Pile domains. The authors kept only large domains because multi-epoch
training violates the assumption of \eqref{eq:preq}.
\item From a random initialisation EpiGen changes the GLUE mean by $-0.003$. No learner improves on
pretrained perplexity.
\item Paper and code disagree. The paper uses 30 domains and 29.7B tokens, the repository 31 and 43.3B.
The EpiGen reward uses a copy of the learner with a fresh AdamW state; each GPU trains a separate
learner; losses are averaged per token instead of summed as in \eqref{eq:epigen}. The paper states four
L40S GPUs for 10 hours per run, the README two GPUs for 24 hours \citep{episelectcode,epigencode}.
\end{enumerate}

% ============================================================================
\section{Citing work and critiques}\label{sec:cite}
% ============================================================================
Semantic Scholar lists 37 papers citing \citet{finzi2026epiplexity} by 23 September 2026
\citep{s2citations}; at least two more exist \citep{ohzeki2026nonequilibrium,okanohara2026thermo}. The
founding paper has no venue; author pages list it as an arXiv preprint. The idea first appeared in
\citet{jiang2025thesis}, which reports the ECA ranking stable across Mamba, temporal-convolution and
transformer observers. \Cref{tab:cite} groups the substantive users.

\begin{table}[!htb]
\centering\footnotesize
\caption{Papers that use epiplexity, grouped by relevance to MLP estimation on small datasets. The
remaining citing papers mention it in one line \citep{s2citations}.}\label{tab:cite}
\begin{tabularx}{\textwidth}{@{}P{1.5cm}P{4.7cm}L@{}}
\toprule
Relevance & Work & Use of epiplexity \\
\midrule
High & \citet{zhang2026learnablenovelty} & closed-form score \eqref{eq:zl}; ECA ranking; objective for automata and an MNIST encoder; intrinsic reward in reinforcement learning, better than task reward on 9 of 10 tasks \\
 & \citet{qiu2026requential} & rigorous estimator \eqref{eq:req} \\
 & \citet{donoway2026edl,donoway2026edlisit} & concurrent EDL \eqref{eq:edl}, compared with $S_T$ \\
 & \citet{li2026counterexample} & proposition and synthetic experiment: a total-structure proxy ranks two pretraining datasets opposite to OOD accuracy (2 of 3 seeds) \\
\midrule
Medium & \citet{liu2026selfplay} & small-data prequential estimate with epoch selection \\
 & \citet{dineen2026vocab} & low-rank adapter as observer, bits per token, epoch chosen by MDL \\
 & \citet{majchrowska2026echo} & ranks self-supervised runs; estimator not described \\
 & \citet{ohzeki2026nonequilibrium} & finite-time information criterion \\
 & \citet{lee2026nca,ren2026creativity} & motivation for automaton pretraining; learner improvement as generation objective \\
 & \citet{takahashi2026thermo} & epiplexity per joule; values benchmark-relative \\
\midrule
Low & \citet{okanohara2026thermo,noguer2026financial,moriondo2026spectral,hu2026speculating,kim2026megadocs,panaganti2026grpo,cornelisse2026autonomy,herrmann2026interestingness} & interpretive or motivating use; no computation \\
\bottomrule
\end{tabularx}
\end{table}

\begin{unsettled}
\textbf{Objections without a resolution.}
\begin{itemize}
\item \emph{Task relevance.} $S_T$ includes structure irrelevant to any downstream task.
\citet{li2026counterexample} constructs datasets whose total-structure proxy ranks opposite to OOD
accuracy; an unreviewed reinforcement-learning curriculum weighted by epiplexity finds no correlation
with transfer
\citep{sunnydigital}. Reason: no task-conditioned variant has been tested.
\item \emph{Looseness of \eqref{eq:preq}.} The founding group calls it a non-rigorous heuristic
\citep{qiu2026requential}. Reason: the ratio to the requential code depends on the teacher recipe.
\item \emph{Finite data.} The one-pass assumption fails on small datasets \citep{liu2026selfplay}.
Reason: the published fixes of \cref{sec:alt} differ by the generalisation gap.
\item \emph{Practicality and relativity.} Each value needs a training sweep, has no gradient with
respect to the data and is benchmark-relative
\citep{zhang2026learnablenovelty,cornelisse2026autonomy,takahashi2026thermo}. Reason: only the
reservoir score is differentiable, with a different observer.
\end{itemize}
\end{unsettled}

% ============================================================================
\section{Theoretical alternatives}\label{sec:theory}
% ============================================================================
\Cref{tab:theory} compares older measures on the criteria that define epiplexity; \cref{tab:practical}
covers code lengths computed from training runs. Measures that stay near zero on random data measure
structure; the others measure randomness.

\begin{table}[!htb]
\centering\footnotesize
\caption{Measures of information and structure. \od: our inference from the definition, not stated by
the source.}\label{tab:theory}
\begin{tabularx}{\textwidth}{@{}P{4.4cm}P{2.6cm}P{1.3cm}P{1.3cm}P{2.2cm}L@{}}
\toprule
Measure & Bounded observer & $\approx0$ on random & $\approx0$ on constant & Estimable & Reference \\
\midrule
Epiplexity $S_T$ & time, training and evaluation & yes & yes & yes & \citet{finzi2026epiplexity} \\
Time-bounded entropy $H_T$ & same & no & yes & held-out loss & \citet{finzi2026epiplexity} \\
Kolmogorov $K$; time-bounded $K^t$, Levin $Kt$ & no; time for $K^t$, $Kt$ & no & yes\od & compressor proxy; exponential time & \citet{livitanyi2019,levin1973,allender2011} \\
Structure function, minimal sufficient statistic; effective complexity & no & yes & yes\od & upper semicomputable at best & \citet{vereshchagin2004,gacs2001,gellmann1996,ay2010} \\
Sophistication (Koppel, naive, coarse) & no; naive time bound collapses & yes & yes\od & no & \citet{koppel1987,antunes2009,mota2013} \\
Logical depth; computational depth & no; they measure time & yes & yes\od & no & \citet{bennett1988,antunes2006} \\
Facticity & no & yes & small & no & \citet{adriaans2009,adriaans2012} \\
Statistical complexity; excess entropy; predictive information & no & yes & yes & block entropies, discrete processes & \citet{shalizi2001,crutchfield2003,bialek2001} \\
HILL, metric, Yao, next-block pseudoentropy & class of tests & no & yes\od & no & \citet{hastad1999,barak2003,haitner2013,vadhan2012} \\
V-entropy & family $\mathcal V$, inference only & no & yes & yes & \citet{xu2020} \\
V-information; pointwise V-information & family $\mathcal V$ & if $X\perp Y$ & yes\od & yes & \citet{xu2020,ethayarajh2022} \\
Two-part MDL, model part & model class & yes\od & yes\od & simple classes & \citet{rissanen1978,grunwald2007} \\
Information in the weights & model class & at $\beta=1$\od & yes\od & Gaussian or Fisher & \citet{achille2019,achille2021} \\
\bottomrule
\end{tabularx}
\end{table}

\paragraph{Relation of each family to $S_T$.} The structure function and the minimal sufficient
statistic split an individual string into model bits and noise bits; $S_T$ is their time-bounded,
distributional analogue (our derivation). Effective complexity is closest: the total information
$K(E)+H(E)$ of an ensemble $E$ is the unbounded form of $|\prog|+\E\log1/P(X)$ \citep{ay2010}.
Sophistication and facticity are uncomputable; naive time-bounded sophistication collapses to a
constant (Lemma~29 of \citealp{finzi2026epiplexity}). Depth measures production time instead of model
size. Excess entropy is the area between finite-block entropy rates and the asymptotic rate, the
unbounded analogue of \eqref{eq:preq} \citep{finzi2026epiplexity}; for such unbounded observers a
pseudorandom stream has large statistical complexity (our derivation). Pseudoentropies are analogues
of $H_T$; next-block pseudoentropy depends on block order, a precedent for Paradox~2 (our
derivation). V-entropy with $\mathcal V=\PT$ is $H_T$ without the model term (our derivation).
V-information measures how much $X$ helps predict $Y$: a simple, highly predictive rule has high
V-information and low $S_T$. Information in the weights \citep{achille2021} is the closest
machine-learning precursor: a divergence to a prior replaces $|\prog|$; no time bound enters.

\begin{remark}[bibliographic corrections]\label{rem:bib}
The founding bibliography lists ``Sophistication revisited'' with four authors; it is by Antunes and
Fortnow alone \citep{antunes2009}, while the four-author paper is ``Computational depth''
\citep{antunes2006}. Coarse sophistication equals busy-beaver depth up to $O(\log n)$ by Theorem~5.3 of
\citet{antunes2009}; \citet{ay2010} relate effective complexity, not sophistication, to logical depth.
Definition~5 of the founding paper omits $x\in S$ (issue~\#1 of \citealp{epiplexitycode}).
\end{remark}

% ============================================================================
\section{Practical neural-network scores}\label{sec:practical}
% ============================================================================
\Cref{tab:practical} sorts practical scores by behaviour on random and trivial labels. Four behave like
$S_T$: the area above a held-out floor; the requential code; surplus description length (SDL) with the
floor at the irreducible loss; per-example learnability. Raw code lengths and difficulty scores rise
with label noise. Consistency and confidence scores are high on trivial data.

\begin{table}[!htb]
\centering\footnotesize
\caption{Practical scores by behaviour on random and trivial labels (our classification of behaviours
reported in the cited papers). $C$: number of classes.}\label{tab:practical}
\begin{tabularx}{\textwidth}{@{}P{5.2cm}P{2.5cm}P{1.6cm}P{1.5cm}P{1.6cm}L@{}}
\toprule
Score & Class & Random labels & Trivial labels & Runs & Reference \\
\midrule
Area above held-out floor; \eqref{eq:preq} in one pass & structure & $\approx0$ & $\approx0$ & 1 & \citet{zhang2020transfer} \\
Requential code & structure & $\approx0$ & $\approx0$ & teacher, student & \citet{qiu2026requential} \\
SDL, floor at irreducible loss & structure & $\approx0$ & $\approx0$ & loss--data curve & \citet{whitney2020} \\
RHO-LOSS, JEST learnability, per example & structure not yet learnt & $\approx0$ or negative & $\approx0$ once learnt & 1 plus reference & \citet{mindermann2022,evans2024} \\
Total prequential code; gzip; bits per character & total information & $\approx M\log C$; high & $\approx0$ & 1 or none & \citet{blier2018,voita2020,pandey2024,deletang2024} \\
Area above training floor after many epochs; intrinsic dimension on training accuracy & capacity incl.\ memorisation & high & low & 1 or many & \citet{zhang2020transfer,li2018intrinsic} \\
EL2N, GraNd, forgetting, gradient variance, margin area & difficulty & high & low & 1--10 & \citet{paul2021,toneva2019,agarwal2022,pleiss2020} \\
C-score; cartography confidence; pointwise V-information & ease & low & high & 1 to thousands & \citet{jiang2021cscore,swayamdipta2020,ethayarajh2022} \\
V-information & predictable information & $\approx0$ & $\approx0$ & 2 & \citet{xu2020} \\
\bottomrule
\end{tabularx}
\end{table}

\eqref{eq:preq} has antecedents the founding paper does not cite. \citet{blier2018} interpret the
prequential code minus the log-loss of the final model as ``the amount of information that the trained
parameters contain about the data''; \citet{voita2020} use the same difference. A text search of v2 finds
neither, nor \citet{bornschein2022}, \citet{perez2021}, \citet{yogatama2019}, RHO-LOSS, pointwise
V-information, EL2N, \citet{achille2021} or \citet{bialek2001}. The paper adopts the held-out heuristic
of \citet{zhang2020transfer}, ``more analogous to the spirit of epiplexity''; it also cites SDL. Variance
reduction transfers directly: incremental training with replay and temperature calibration
\citep{bornschein2022}; switching between models against early-block overfitting \citep{blier2018};
several seeds \citep{perez2021}.

% ============================================================================
\section{Consequences for small models on one GPU}\label{sec:small}
% ============================================================================
No published work estimates epiplexity with an MLP or a convolutional network; every trained
experiment of the founding group uses transformers
\citep{finzi2026epiplexity,qiu2026requential,su2026epiplexity}. Several ingredients already exist at
small scale (\cref{tab:small}).

\begin{table}[!htb]
\centering\footnotesize
\caption{The smallest published or measured settings relevant to MLP estimation.}\label{tab:small}
\begin{tabularx}{\textwidth}{@{}P{5.4cm}P{3.0cm}P{2.6cm}L@{}}
\toprule
Setting & Size & Compute or time & Value \\
\midrule
ECA frontier, smallest models \citep{finzi2026epiplexity} & $\approx3\times10^3$ parameters & order settles from $\approx10^{13}$ FLOPs & $S_T\approx10^5$--$10^7$ bits \\
Prequential code of MNIST, CIFAR-10 labels \citep{blier2018} & $6\times10^4$, $5\times10^4$ labels & not reported & 4.10, 45.3 kbits; uniform code 199, 166 \\
MNIST LeNet with replay \citep{bornschein2022} & $6\times10^4$ labels & $1.4\times10^{13}$ FLOPs & $4.4\times10^3$ nats \\
MLP classifiers on frozen MNIST features \citep{whitney2020} & 20 sizes, 8 seeds & $\approx2$\,min, one GPU & loss--data curves \\
Reservoir score, one ECA rule\od & $N_{\rm s}=512$ & $\approx0.7$\,s, CPU & \cref{tab:zl} \\
Released PyTorch trainer, rule 54, 1 layer, width 32\od & 200 steps & 143\,s, CPU & area $\approx8.3\times10^5$ nats \\
\bottomrule
\end{tabularx}
\end{table}

\paragraph{Pitfalls.} Seven problems decide whether a small-scale value means anything.
\begin{enumerate}
\item \emph{Training-set floor.} After several epochs a training-set floor turns memorised noise into
apparent structure. Random labels on the $M=6\times10^4$ MNIST images would then score up to
$M\log_2 10\approx199$ kbits instead of zero (our derivation). A held-out floor keeps \eqref{eq:preq}
near zero \citep{zhang2020transfer}; EDL makes this a theorem.
\item \emph{Floor error times $M$.} A bias $\delta$ in the final loss shifts \eqref{eq:preq} by
$M\delta$. At the ECA scale, $\delta=10^{-3}$ nats over $2.5\times10^8$ tokens gives $3.6\times10^5$
bits, more than $S_T$ of rule 15 (our derivation).
\item \emph{Magnitudes.} The MNIST prequential code bounds the area above any non-negative floor
(\cref{tab:small}); \cref{prop:pl}(c) scales $S_\infty$ as $\Dt^{1-\beta}$. Small datasets give
kilobits, against megabits for the synthetic tasks of \cref{tab:exp} (our derivation).
\item \emph{Requential overhead.} At batch 128 one MNIST epoch has 469 steps; the constant of
\eqref{eq:req} alone adds $469\kappa\approx2.4$ kbits, the order of the whole MNIST prequential code
(our derivation).
\item \emph{Units and token bookkeeping.} The repositories log nats; Fig.~2c of the founding paper is
labelled ``MB'' and the hard-induction notebook plots nats. The released prequential area integrates
over all tokens while the loss covers the label half; the requential sum covers label tokens only.
Part of the gap in Fig.~2c may be this factor of two (our derivation from the code).
\item \emph{No seeds.} No epiplexity value in \cref{sec:exp,sec:follow} has an error bar. Related MDL
studies use four to eight seeds \citep{voita2020,perez2021,whitney2020}.
\item \emph{One compute value.} Fig.~3 of the founding paper ranks rule 15 above rule 54 at the lowest
compute. Plateau-then-drop tasks such as sparse parity make the area depend on seed and truncation
\citep{barak2022parity}.
\end{enumerate}

\begin{stated}
The literature fixes five requirements for a small-scale estimate. The floor of \eqref{eq:preq} is a
held-out or first-pass loss. The standard error of the floor is small against $S/M$. Code lengths are
in bits over label units, with $\Dt$ and $T$ stated. Requential values at small batch include the
per-message terms of \eqref{eq:req}. Rankings are checked at more than one compute value.
\end{stated}

\paragraph{Testbeds.} \Cref{tab:testbeds} ranks candidate datasets by a variable that sets the amount
of structure, a known entropy floor and known learnability. Two designs isolate single effects (our
derivation). Symmetric label noise at rate
$\rho$ over $C$ classes sets $H(Y\mid X)=h(\rho)+\rho\log_2(C-1)$ per example, with $h$ the binary
entropy; the raw prequential code grows with this term, while \eqref{eq:preqkl} and the requential
code should not. A fixed pixel permutation is an exact symmetry of an MLP with i.i.d.\ initialisation
but destroys locality for a convolutional network. $S_T$ should be unchanged for the first observer
and change for the second.

\begin{table}[!htb]
\centering\footnotesize
\caption{Candidate testbeds, ranked by our assessment of fitness for MLP estimation on one GPU.}\label{tab:testbeds}
\begin{tabularx}{\textwidth}{@{}rP{4.0cm}P{2.3cm}P{2.1cm}LP{2.5cm}@{}}
\toprule
& Testbed & Known floor & Structure variable & Main risk & Source \\
\midrule
1 & ECA, $Y=F^t(X)$ & zero & rule; steps $t$ & an MLP lacks locality; add a 1-D convolutional network & \citet{finzi2026epiplexity} \\
2 & Markov sources; hidden-row induction & entropy rate & hidden rows & induction needs a sequence model & \citet{edelman2024induction} \\
3 & true, noisy, random labels; permuted pixels (MNIST, CIFAR-10, MNIST-1D) & $h(\rho)+\rho\log_2(C-1)$\od & noise rate; permutation & kilobit signal & \citet{zhang2017rethinking,greydanus2020} \\
4 & random hierarchy model & sample complexity known & branching; depth & custom generator & \citet{cagnetta2024rhm} \\
5 & sparse parity $(n,k)$ & zero & $k$ & plateau then drop & \citet{barak2022parity} \\
6 & modular arithmetic & zero & modulus & finite table forces many epochs & \citet{nanda2023,gromov2023} \\
7 & Game of Life & zero & steps & needs strong overparameterisation & \citet{springer2020life} \\
8 & linear congruential against cryptographic generators & $S_T\approx0$ for a PRG & generator & transformers needed A100-hours & \citet{tao2025lcg} \\
\bottomrule
\end{tabularx}
\end{table}

\paragraph{Code.} The PyTorch trainer of \citet{epiplexitycode} runs on a CPU as a library call with
logging disabled and returns the prequential area per checkpoint (our measurement). An MLP port needs a
model class with the same interface, per-class sampling in place of autoregressive generation and
areas over label units. A full reproduction of Fig.~3 would take two to four T4-days, a large share of
it in data generation on the CPU; a reduced grid (widths 16--128, depths 1--2, data generated on the
GPU) takes under two hours per rule (our estimate). \Cref{tab:repos} lists the reusable pieces.

\paragraph{Protocol implied by the literature} (our design).
\begin{enumerate}
\item Fix $\Dt$ and three to five values of $T$ for all datasets; report bits, bits per label and
$S_T/\MDL_T$.
\item Use one pass over fresh data where the generator allows; otherwise use the first-epoch online
loss with a held-out floor, as in \eqref{eq:edl}.
\item Sweep four to six widths with learning rates scaled by $1/\text{fan-in}$, a constant rate and an
EMA; run requential coding with the exact class divergence and report \eqref{eq:req} beside
$\sum_i\KL_i$.
\item Use at least five seeds with a paired bootstrap; include random labels, constant labels and
shuffled targets at matched size.
\end{enumerate}

\begin{unsettled}
\textbf{Questions an MLP project on one GPU can answer.}
\begin{itemize}
\item \emph{Does an MLP estimate of $S_T$ reproduce the ECA ordering?} Reason: no published estimate
uses an MLP or a convolutional network.
\item \emph{Does $S_T$ depend on order at small scale?} Reason: only chess shows it.
\item \emph{How large is the seed spread of $S_T$ and of the hull?} Reason: no source reports it.
\item \emph{Does requential coding work for classifiers?} Reason: no source validates it; at small
batch the per-message terms of \eqref{eq:req} are the order of the signal.
\item \emph{Which teacher recipe is better, a divergence threshold or iso-loss projection?} Reason:
the two papers never compare them.
\item \emph{Do the fast estimators agree with $S_T$?} Reason: EDL, the reservoir score and the
scaling-law proxy have never been applied to the data of a compute sweep.
\item \emph{Does $S_T$ predict transfer for small models?} Reason: the evidence is five Pile domains,
one counterexample and one null result.
\end{itemize}
\end{unsettled}

\FloatBarrier
\bibliographystyle{plainnat}
\bibliography{refs}

\appendix
% ============================================================================
\section{Repositories and discrepancies}\label{sec:repos}
% ============================================================================
\begin{table}[!htb]
\centering\footnotesize
\caption{Repositories. Dates are the last commit, or creation where only that is known.}\label{tab:repos}
\begin{tabularx}{\textwidth}{@{}P{3.4cm}P{2.3cm}P{1.8cm}P{1.8cm}L@{}}
\toprule
Repository & Framework & Licence & Date & Reusable for small scale \\
\midrule
\citet{epiplexitycode} & PyTorch; JAX for natural data & MIT & 8 Mar 2026 & \texttt{soph/train.py}: prequential area, requential loop with EMA teacher and threshold; hull code in \texttt{notebooks/eca\_3rules.ipynb} \\
\citet{requentialcode} & JAX, Flax, Hydra & MIT & 14 Jul 2026 & \texttt{train.py} with iso-loss projection; noise and repeat controls need no data \\
\citet{episelectcode} & JAX, Equinox, jaxopt & MIT & 25 Aug 2026 & \texttt{scaling\_law.py}: joint fit and gradient \eqref{eq:episelect}; per-domain area in \texttt{ado.py}; the fit bounds $\log\epsilon\ge0.5$, too high for low-loss classifiers\od \\
\citet{epigencode} & PyTorch, Hugging Face & MIT & 13 Aug 2026 & look-ahead reward \texttt{\_loss\_drop} \\
\citet{learnablenoveltycode} & PyTorch & Apache 2.0, non-commercial & not recorded & \texttt{core.py}: score \eqref{eq:zl}; \texttt{online.py}: recursive least squares \\
\citet{selfplaycode} & PyTorch, PEFT & not stated & 7 Jul 2026 & \texttt{epiplexity\_mdl.py}: epoch-selection protocol \\
\citet{bossmanlab} & numpy & not stated & 24 May 2026 & next-token MLP transfer experiment \\
\citet{edirentaudio} & PyTorch & not stated & 16 Jun 2026 & synthetic controls: silence, tones, noise \\
\citet{lcrhepiplexity} & TensorFlow.js & not stated & not recorded & reservoir and convolutional observers in the browser \\
\citet{lmxxfprng} & PyTorch & not stated & not recorded & learnability boundary of pseudorandom generators \\
\citet{richardscottoz,taiyou,tunedorg} & Python & not stated & from Jan 2026 & estimator wrappers without published validation \\
\bottomrule
\end{tabularx}
\end{table}

Discrepancies between papers and code, from a comparison of both:
\begin{itemize}
\item \texttt{eca\_3rules.py} uses batch 384 and tag \texttt{arxiv\_3rules}; App.~C.1 states batch
1536 and the notebook filters tag \texttt{test} \citep{epiplexitycode}.
\item \texttt{eca\_rules.py} disables requential coding and runs 128 automaton steps; Fig.~2c plots
requential values and App.~C.7 states 48 steps.
\item \texttt{induction\_easy.py} sets a divergence threshold of 0.005; the paper states none.
\item The caption of Fig.~8a states $10^9$ tokens and the text $5\times10^9$. The totals imply about
1.4 bits per character only for $5\times10^9$ (our derivation).
\item Theorem~18 in the appendix keeps the v1 wording ``CSPRNG'' and ``$-O(1)$''.
\item The requential paper gives the teacher EMA window as $\max(50,at)$ next to $a=0.01$; the
repository sets $a=0.1$ for the teacher \citep{requentialcode}.
\item App.~C.1 of \citet{su2026epiplexity} names Pile domains in a Common Pile analysis. The EpiGen
mixture configurations replace 25, 50 and 100 per cent of learner data by real text, while the paper
describes 25 and 50 per cent synthetic \citep{epigencode}.
\item App.~D of \citet{zhang2026learnablenovelty} places rule 41 second, Fig.~2 rule 25. The README
lists a table column absent from the paper.
\item The arXiv EDL preprint credits one diagnostic example with ``exactly'' $\log|\mathcal Y|$ bits
for label alphabet $\mathcal Y$ in Sec.~5.2 and ``up to'' $\log m$ bits for $m$ hypotheses in
App.~B.2 \citep{donoway2026edl}.
\end{itemize}

\end{document}
